\begin{prerequis}
Connaître le théorème de Pythagore, les rapports trigonométriques dans un triangle
rectangle et les angles usuels en degrés.
\end{prerequis}
\begin{automatismes}
Convertir un demi-tour, un quart de tour et un sixième de tour en degrés; placer
ces angles sur un cercle; rappeler $\sin 30^\circ$ et $\cos 60^\circ$.
\end{automatismes}
\begin{activite}[Activité d'entrée -- Mesurer un angle par une longueur]
Sur des cercles de rayons différents, construire un arc de longueur égale au
rayon. Comparer l'angle au centre obtenu. Cette invariance conduit à définir le
radian; en déduire la mesure d'un tour complet.
\end{activite}

